% Slide: Bối cảnh đề tài
\begin{frame}{Bối Cảnh Đề Tài: Chăm Sóc Thú Cưng Đô Thị}
  \framesubtitle{\color{PetcareSecondary}\textbf{Xu hướng nhân hóa thú cưng \& Nghịch lý thời gian của chủ nuôi bận rộn}}

  \begin{columns}[T]
    % Cột 1: 3 Chuyển dịch đô thị then chốt
    \begin{column}{0.54\textwidth}
      \begin{block}{\footnotesize \textbf{Sự chuyển dịch văn hóa \& Lối sống đô thị}}
        \begin{itemize}
          \setlength\itemsep{1.2mm}
          \tiny
          \item \textbf{Hiện tượng nhân hóa thú cưng (Pet Humanization):} Thú cưng trở thành thành viên ruột thịt trong gia đình; nhu cầu spa, cắt tỉa và chăm sóc sức khỏe định kỳ tăng trưởng bùng nổ tại TP.HCM.
          \item \textbf{Chân dung chủ nuôi bận rộn:} Đa số là nhân viên văn phòng, chuyên viên, người trẻ có lịch làm việc dày đặc, eo hẹp quỹ thời gian trong tuần.
          \item \textbf{Nhu cầu an toàn cao độ:} Thú cưng có vị trí tình cảm đặc biệt, chủ nuôi luôn lo lắng cao độ khi gửi bé tại cơ sở bên ngoài nhưng thiếu công cụ tương tác số hóa.
        \end{itemize}
      \end{block}
    \end{column}

    % Cột 2: Thực trạng quy trình thủ công phân mảnh
    \begin{column}{0.44\textwidth}
      \begin{exampleblock}{\footnotesize \textbf{Thực trạng vận hành cơ sở spa hiện nay}}
        \tiny
        \begin{itemize}
          \setlength\itemsep{1mm}
          \item \textbf{Kênh gọi điện hotline:} Vừa nghe máy vừa lật sổ tay đối soát; dễ nhầm lẫn hoặc trùng lịch.
          \item \textbf{Kênh chat Facebook/Zalo:} Tin nhắn trôi dạt, nhân viên trả lời chậm từ 15 phút đến vài tiếng.
          \item \textbf{Ghi chép tiếp nhận tại chỗ:} Sổ sách giấy, dán giấy post-it lên lồng dễ rơi mất khi đổi ca trực.
        \end{itemize}
      \end{exampleblock}
    \end{column}
  \end{columns}

  \vspace{1.5mm}
  \begin{alertblock}{\centering \scriptsize \textbf{Nghịch lý cốt lõi}}
    \centering \tiny
    Nhu cầu chăm sóc chuyên nghiệp tăng cao nhưng quy trình tương tác vẫn kẹt lại ở thời kỳ \textbf{thủ công, phân mảnh và phụ thuộc hoàn toàn vào trí nhớ cá nhân}.
  \end{alertblock}
\end{frame}
